\chapter{Vectoren}\label{ch:g10:vectors}

Een \omterm{def:g10:vectors:vector}{vector} legt een verplaatsing vast: een richting en een lengte, los van
het beginpunt. \omterm{def:g10:vectors:vector}{Vectoren} leveren nette bewijzen van meetkundige feiten en
herleiden ze, zodra er \omterm{def:g10:coordgeom:system}{coördinaten} bij komen, tot kleine berekeningen. Ze
worden een sleutelgereedschap in \cref{ch:g10:lines} en, later, voor het
scalair product in \cref{ch:g11:scal}.

\section{Verschuivingen en vectoren}

\begin{definition}[Vector]\label{def:g10:vectors:vector}
Bij twee punten $A$ en $B$ stelt de \emph{vector}\index{vector} $\vect{AB}$
de verplaatsing van $A$ naar $B$ voor: hij heeft een \emph{richting} (die van
de rechte $(AB)$, samen met de zin waarin je haar doorloopt, van $A$ naar
$B$) en een \emph{lengte} (de afstand $AB$, ook geschreven als
$\norm{\vect{AB}}$). Twee vectoren zijn \emph{gelijk} wanneer ze dezelfde
verplaatsing vastleggen: $\vect{AB} = \vect{CD}$ betekent dat de twee
verplaatsingen op elk punt hetzelfde resultaat geven.
\end{definition}

\begin{proposition}[Gelijke vectoren en parallellogrammen]\label{prop:g10:vectors:parallelogram}
$\vect{AB} = \vect{CD}$ geldt precies wanneer $ABDC$ (in die volgorde!) een
parallellogram is, eventueel een platgedrukt --- of gelijkwaardig: wanneer
$[AD]$ en $[BC]$ hetzelfde \omterm{prop:g10:coordgeom:midpoint}{midden} hebben.
\end{proposition}
\admitted

\begin{omfigure}
\begin{tikzpicture}[scale=1.0]
  \coordinate (A) at (0,0);
  \coordinate (B) at (2.2,0.8);
  \coordinate (C) at (1.0,-1.2);
  \coordinate (D) at (3.2,-0.4);
  \draw[omExo, thin] (A) -- (C);
  \draw[omExo, thin] (B) -- (D);
  \draw[-Stealth, omDef, thick] (A) -- (B)
    node[midway, above, font=\small] {$\vect{AB}$};
  \draw[-Stealth, omDef, thick] (C) -- (D)
    node[midway, below, font=\small] {$\vect{CD}$};
  \fill (A) circle (1.5pt) node[left, font=\small] {$A$};
  \fill (B) circle (1.5pt) node[right, font=\small] {$B$};
  \fill (C) circle (1.5pt) node[left, font=\small] {$C$};
  \fill (D) circle (1.5pt) node[right, font=\small] {$D$};
\end{tikzpicture}

{\small Gelijke \omterm{def:g10:vectors:vector}{vectoren} $\vect{AB} = \vect{CD}$: dezelfde richting, dezelfde
zin, dezelfde lengte. De vierhoek $ABDC$ is een parallellogram.}
\end{omfigure}

\begin{notation}
Een \omterm{def:g10:vectors:vector}{vector} krijgt vaak één letter als naam, $\vec u$ of $\vec v$, wanneer de
eindpunten er niet toe doen. De \emph{nulvector} $\vec 0 = \vect{AA}$ is de
verplaatsing die niets verplaatst. De \emph{tegengestelde} van $\vect{AB}$ is
$\vect{BA}$: dezelfde lengte, tegengestelde zin. We schrijven
$\vect{BA} = -\vect{AB}$.
\end{notation}

\section{Vectoren optellen}

\begin{definition}[Som van twee vectoren]\label{def:g10:vectors:sum}
De \emph{som}\index{vector!som} $\vec u + \vec v$ is de verplaatsing die je
krijgt door eerst $\vec u$ en daarna $\vec v$ uit te voeren.
\end{definition}

\begin{theorem}[Relatie van Chasles]\label{thm:g10:vectors:chasles}
Voor elke drie punten $A$, $B$, $C$ geldt
\[
\vect{AB} + \vect{BC} = \vect{AC}.
\]
\end{theorem}

\begin{proof}
Van $A$ naar $B$ gaan en daarna van $B$ naar $C$ is een verplaatsing die $A$
naar $C$ brengt: precies de verplaatsing die $\vect{AC}$ vastlegt.
\end{proof}

\begin{remark}[Parallellogramregel]
Om twee \omterm{def:g10:vectors:vector}{vectoren} op te tellen die vanuit hetzelfde punt vertrekken,
$\vect{AB} + \vect{AC}$, vervolledig je het parallellogram $ABDC$: de \omterm{def:g10:vectors:sum}{som} is
de diagonaal $\vect{AD}$. Inderdaad is $\vect{AC} = \vect{BD}$, dus
$\vect{AB} + \vect{AC} = \vect{AB} + \vect{BD} = \vect{AD}$ volgens de
relatie van Chasles.
\end{remark}

\begin{omfigure}
\begin{tikzpicture}[scale=1.05]
  \coordinate (A) at (0,0);
  \coordinate (B) at (2.4,0.5);
  \coordinate (C) at (3.4,2.0);
  \draw[-Stealth, omDef, thick] (A) -- (B)
    node[midway, below, font=\small] {$\vec u$};
  \draw[-Stealth, omThm, thick] (B) -- (C)
    node[midway, right, font=\small] {$\vec v$};
  \draw[-Stealth, omMeth, thick] (A) -- (C)
    node[midway, above left, font=\small] {$\vec u + \vec v$};
\end{tikzpicture}\hspace{1.4cm}%
\begin{tikzpicture}[scale=1.05]
  \coordinate (A) at (0,0);
  \coordinate (B) at (2.4,0.5);
  \coordinate (C) at (1.0,1.5);
  \coordinate (D) at (3.4,2.0);
  \draw[omExo, thin] (B) -- (D) -- (C);
  \draw[-Stealth, omDef, thick] (A) -- (B)
    node[midway, below, font=\small] {$\vect{AB}$};
  \draw[-Stealth, omThm, thick] (A) -- (C)
    node[midway, left, font=\small] {$\vect{AC}$};
  \draw[-Stealth, omMeth, thick] (A) -- (D)
    node[pos=0.75, below right, font=\small] {$\vect{AD}$};
  \fill (D) circle (1.2pt) node[right, font=\small] {$D$};
\end{tikzpicture}

{\small Twee \omterm{def:g10:functions:function}{beelden} van de \omterm{def:g10:vectors:sum}{som}: kop aan staart (relatie van Chasles, links)
en de parallellogramregel (rechts).}
\end{omfigure}

\begin{example}[Vereenvoudigen met Chasles]\label{ex:g10:vectors:chasles}
Vereenvoudig $\vect{MN} + \vect{NP} + \vect{PQ}$: de verplaatsingen aan
elkaar rijgen geeft $\vect{MQ}$. Vereenvoudig $\vect{AB} - \vect{AC}$:
herschrijf het verschil als een \omterm{def:g10:vectors:sum}{som} met de tegengestelde \omterm{def:g10:vectors:vector}{vector},
\[
\vect{AB} - \vect{AC} = \vect{AB} + \vect{CA}
= \vect{CA} + \vect{AB} = \vect{CB}.
\]
\end{example}

\section{Een vector met een getal vermenigvuldigen}

\begin{definition}[Scalair veelvoud]\label{def:g10:vectors:scalar}
Zij $\vec u$ een \omterm{def:g10:vectors:vector}{vector} verschillend van $\vec 0$ en $k$ een \omterm{def:g10:numbers:sets}{reëel getal}. De
\omterm{def:g10:vectors:vector}{vector} $k\vec u$\index{vector!scalair veelvoud} heeft dezelfde richting als
$\vec u$, lengte $\abs{k} \times \norm{\vec u}$, en dezelfde zin als
$\vec u$ wanneer $k > 0$, de tegengestelde zin wanneer $k < 0$. Voor $k = 0$
is $0\vec u = \vec 0$.
\end{definition}

\begin{definition}[Collineariteit]\label{def:g10:vectors:collinear}
Twee \omterm{def:g10:vectors:vector}{vectoren} $\vec u$ en $\vec v$ heten
\emph{collineair}\index{collineaire vectoren} wanneer ze dezelfde richting
hebben, dus wanneer $\vec v = k \vec u$ voor een zeker \omterm{def:g10:numbers:sets}{reëel getal} $k$ (of
wanneer een van beide $\vec 0$ is).
\end{definition}

\begin{remark}
Collineariteit is de vectortaal voor evenwijdigheid en voor het op één rechte
liggen:
\begin{itemize}
  \item de rechten $(AB)$ en $(CD)$ zijn evenwijdig precies wanneer
        $\vect{AB}$ en $\vect{CD}$ \omterm{def:g10:vectors:collinear}{collineair} zijn;
  \item de punten $A$, $B$, $C$ liggen op één rechte precies wanneer
        $\vect{AB}$ en $\vect{AC}$ \omterm{def:g10:vectors:collinear}{collineair} zijn.
\end{itemize}
\end{remark}

\section{Coördinaten van vectoren}

\begin{definition}[Coördinaten van een vector]\label{def:g10:vectors:coords}
In een \omterm{def:g10:coordgeom:system}{assenstelsel} zijn de \emph{\omterm{def:g10:coordgeom:system}{coördinaten}} van de \omterm{def:g10:vectors:vector}{vector} $\vect{AB}$ de
getallen die de verplaatsing beschrijven:
\[
\vect{AB}\,(x_B - x_A,\ y_B - y_A).
\]
Een \omterm{def:g10:vectors:vector}{vector} $\vec u\,(a, b)$ verplaatst elk punt over $a$ eenheden horizontaal
en $b$ eenheden verticaal.
\end{definition}

\begin{proposition}[Rekenen met coördinaten]\label{prop:g10:vectors:coordrules}
Zij $\vec u\,(a, b)$ en $\vec v\,(c, d)$, en $k \in \R$.
\begin{enumerate}
  \item $\vec u = \vec v$ precies wanneer $a = c$ en $b = d$;
  \item $\vec u + \vec v$ heeft \omterm{def:g10:coordgeom:system}{coördinaten} $(a + c,\ b + d)$;
  \item $k\vec u$ heeft \omterm{def:g10:coordgeom:system}{coördinaten} $(ka,\ kb)$;
  \item $\norm{\vec u} = \sqrt{a^2 + b^2}$ (\omterm{def:g10:coordgeom:system}{orthonormaal assenstelsel}).
\end{enumerate}
\end{proposition}

\begin{proof}
1--3 herhalen, coördinaat per coördinaat, wat de verplaatsingen doen: eerst
$\vec u$ en dan $\vec v$ uitvoeren verplaatst een punt horizontaal eerst over
$a$ en dan over $c$, dus in totaal over $a + c$. Punt 4 is de afstandsformule
van \cref{thm:g10:coordgeom:distance}, toegepast op een lijnstuk dat
$\vec u$ voorstelt.
\end{proof}

\begin{theorem}[Criterium voor collineariteit]\label{thm:g10:vectors:det}
Twee \omterm{def:g10:vectors:vector}{vectoren} $\vec u\,(a, b)$ en $\vec v\,(c, d)$ zijn \omterm{def:g10:vectors:collinear}{collineair} dan en
slechts dan als
\[
ad - bc = 0 .
\]
\end{theorem}

\begin{proof}
Is $\vec v = k\vec u$, dan is $c = ka$ en $d = kb$, dus
$ad - bc = a(kb) - b(ka) = 0$. Hetzelfde geldt als $\vec u = k \vec v$ of als
een van beide \omterm{def:g10:vectors:vector}{vectoren} nul is.

Omgekeerd, stel $ad - bc = 0$ met $\vec u \neq \vec 0$, zeg $a \neq 0$ (het
geval $b \neq 0$ verloopt net zo). Stel $k = \frac{c}{a}$; dan is $c = ka$,
en uit $ad = bc = b(ka)$ volgt na deling door $a \neq 0$ dat $d = kb$. Dus is
$\vec v = k\vec u$.
\end{proof}

\begin{example}\label{ex:g10:vectors:det}
Zijn $\vec u\,(3, -2)$ en $\vec v\,(-6, 4)$ \omterm{def:g10:vectors:collinear}{collineair}? Bereken
$ad - bc = 3 \times 4 - (-2)\times(-6) = 12 - 12 = 0$: ja, en inderdaad is
$\vec v = -2\vec u$. Voor $\vec u\,(3, -2)$ en $\vec w\,(1, 5)$:
$3 \times 5 - (-2) \times 1 = 17 \neq 0$: niet \omterm{def:g10:vectors:collinear}{collineair}.
\end{example}

\begin{method}[Op één rechte liggen en evenwijdigheid]\label{met:g10:vectors:alignment}
Om te bewijzen dat drie punten $A$, $B$, $C$ op één rechte liggen:
\begin{enumerate}
  \item bereken de \omterm{def:g10:coordgeom:system}{coördinaten} van $\vect{AB}$ en $\vect{AC}$;
  \item ga het criterium $ad - bc = 0$ van \cref{thm:g10:vectors:det} na;
  \item besluit: de \omterm{def:g10:vectors:vector}{vectoren} zijn \omterm{def:g10:vectors:collinear}{collineair} en delen het punt $A$, dus
        liggen de drie punten op één rechte.
\end{enumerate}
Dezelfde berekening met $\vect{AB}$ en $\vect{CD}$ bewijst dat de rechten
$(AB)$ en $(CD)$ evenwijdig zijn.
\end{method}

\begin{example}\label{ex:g10:vectors:alignment}
Zij $A(1, 2)$, $B(3, 5)$ en $C(7, 11)$. Dan is $\vect{AB}\,(2, 3)$ en
$\vect{AC}\,(6, 9)$, en $2 \times 9 - 3 \times 6 = 18 - 18 = 0$: de punten
liggen op één rechte (er geldt zelfs $\vect{AC} = 3\vect{AB}$).
\end{example}

\begin{omfigure}
\begin{tikzpicture}[scale=0.62]
  \draw[thin, gray!40] (-0.4,-0.4) grid (8.4,5.4);
  \draw[-Stealth] (-0.6,0) -- (8.6,0);
  \draw[-Stealth] (0,-0.6) -- (0,5.6);
  \coordinate (A) at (1,1);
  \coordinate (B) at (3,2);
  \coordinate (C) at (7,4);
  \draw[-Stealth, omDef, very thick] (A) -- (B)
    node[midway, above left=-2pt, font=\small] {$\vect{AB}$};
  \draw[-Stealth, omThm, thick] (A) -- (C)
    node[pos=0.8, below right=-2pt, font=\small] {$\vect{AC} = 3\vect{AB}$};
  \fill (A) circle (2pt) node[above left, font=\small] {$A$};
  \fill (B) circle (2pt) node[above left, font=\small] {$B$};
  \fill (C) circle (2pt) node[above left, font=\small] {$C$};
\end{tikzpicture}

{\small Punten op één rechte: $\vect{AC}$ is een scalair veelvoud van
$\vect{AB}$.}
\end{omfigure}

\section{Oefeningen}

\begin{exercise}[$\star$]\label{exo:g10:vectors:1}
Vereenvoudig met de relatie van Chasles:
\[
\vect{AB} + \vect{BD}, \qquad
\vect{KL} + \vect{LM} + \vect{MK}, \qquad
\vect{AC} - \vect{BC}, \qquad
\vect{AB} + \vect{CA}.
\]
\end{exercise}

\begin{exercise}[$\star$]\label{exo:g10:vectors:2}
Zij $A(2, 1)$, $B(5, 3)$, $C(-1, 4)$. Bereken de \omterm{def:g10:coordgeom:system}{coördinaten} van
$\vect{AB}$, $\vect{BC}$, $\vect{AC}$, en ga op de \omterm{def:g10:coordgeom:system}{coördinaten} na dat
$\vect{AB} + \vect{BC} = \vect{AC}$.
\end{exercise}

\begin{exercise}[$\star$]\label{exo:g10:vectors:3}
Zij $\vec u\,(2, -3)$ en $\vec v\,(-1, 4)$. Geef de \omterm{def:g10:coordgeom:system}{coördinaten} van
$\vec u + \vec v$, $3\vec u$, $2\vec u - \vec v$, en bereken
$\norm{\vec u}$.
\end{exercise}

\begin{exercise}[$\star$]\label{exo:g10:vectors:4}
Zeg in elk geval of $\vec u$ en $\vec v$ \omterm{def:g10:vectors:collinear}{collineair} zijn:
\[
\text{(a) } \vec u\,(4, 6),\ \vec v\,(6, 9);
\qquad
\text{(b) } \vec u\,(2, -5),\ \vec v\,(-4, 10);
\qquad
\text{(c) } \vec u\,(3, 1),\ \vec v\,(1, 3).
\]
\end{exercise}

\begin{exercise}[$\star$]\label{exo:g10:vectors:5}
Zij $A(0, 2)$, $B(4, 0)$, $C(6, 3)$. Zoek de \omterm{def:g10:coordgeom:system}{coördinaten} van het punt $D$
waarvoor $ABCD$ een parallellogram is, dus waarvoor
$\vect{AB} = \vect{DC}$.
\end{exercise}

\begin{exercise}[$\star\star$]\label{exo:g10:vectors:6}
Zij $A(-2, 1)$, $B(1, 3)$, $C(4, -1)$.
\begin{enumerate}
  \item Bereken de \omterm{def:g10:coordgeom:system}{coördinaten} van het punt $M$ waarvoor
        $\vect{AM} = \vect{AB} + \vect{AC}$.
  \item Wat is de vierhoek $ABMC$? Verantwoord.
\end{enumerate}
\end{exercise}

\begin{exercise}[$\star\star$]\label{exo:g10:vectors:7}
Zij $A(1, -1)$, $B(4, 1)$, $C(10, 5)$. Liggen de punten $A$, $B$, $C$ op één
rechte? Dezelfde vraag voor $A(0, 3)$, $B(2, 2)$, $C(8, -1)$.
\end{exercise}

\begin{exercise}[$\star\star$]\label{exo:g10:vectors:8}
Zij $A(1, 2)$, $B(5, 4)$, $C(6, 1)$, $D(2, -1)$. Toon aan dat $(AB)$ en
$(CD)$ evenwijdig zijn, en daarna dat $ABCD$ een parallellogram is. Welke
vaststelling houdt de andere in?
\end{exercise}

\begin{exercise}[$\star\star$]\label{exo:g10:vectors:9}
Zij $ABC$ een driehoek en $I$ het punt bepaald door
$\vect{AI} = \frac23 \vect{AB}$. Neem $A(0,0)$, $B(6, 3)$, $C(2, 5)$ (zodat
de berekeningen eenvoudig blijven):
\begin{enumerate}
  \item bereken de \omterm{def:g10:coordgeom:system}{coördinaten} van $I$;
  \item zij $J$ zo dat $\vect{CJ} = \frac23\vect{CB}$; bereken de \omterm{def:g10:coordgeom:system}{coördinaten}
        van $J$;
  \item toon aan dat $(IJ)$ evenwijdig is met $(AC)$.
\end{enumerate}
\end{exercise}

\begin{exercise}[$\star\star\star$]\label{exo:g10:vectors:10}
Zij $ABCD$ een willekeurige vierhoek, en zijn $I$, $J$, $K$, $L$ de \omterm{prop:g10:coordgeom:midpoint}{middens}
van $[AB]$, $[BC]$, $[CD]$, $[DA]$. Toon met \omterm{def:g10:coordgeom:system}{coördinaten} $A(x_A, y_A)$
enzovoort aan dat $\vect{IJ} = \vect{LK}$, en besluit dat de \omterm{prop:g10:coordgeom:midpoint}{middens} van de
zijden van \emph{elke} vierhoek een parallellogram vormen.
\end{exercise}

\section{Opgave: de algebra van de pijlen}

\begin{problem}[{Weekendopgave --- gymnastiek met Chasles, het zwaartepunt
in twee regels opnieuw bewezen, en krachten in evenwicht: waar vectoren
werkelijk voor dienen}]
\label{pb:g10:vectors:1}
In jaar 8 werd bewezen dat de zwaartelijnen van een driehoek elkaar op twee
derde van hun lengte ontmoeten --- met een vernuftig parallellogram verborgen
in de driehoek (de weekendopgave over het zwaartepunt in het
onderbouwvolume). \omterm{def:g10:vectors:vector}{Vectoren} bewijzen het in twee regels, en geven de
concurrentie er gratis bij. Daar dient deze nieuwe algebra voor: stellingen
worden berekeningen met pijlen. Deze opgave oefent het rekenwerk (Chasles
vooral), levert het bewijs van twee regels, breidt Varignon uit
(\cref{exo:g10:vectors:10}) en eindigt waar de \omterm{def:g10:vectors:vector}{vectoren} geboren zijn:
krachten en snelheden.

\medskip
\noindent\textbf{Deel I --- Gymnastiek met Chasles.}
\begin{enumerate}
  \item Vereenvoudig met de relatie van Chasles
        (\cref{thm:g10:vectors:chasles}):
        $\vect{AB} + \vect{BC} + \vect{CD}$; \quad
        $\vect{MA} - \vect{MB}$; \quad
        $\vect{AB} + \vect{CD} + \vect{BC} + \vect{DA}$.
  \item De oorsprongtruc: toon aan dat voor \emph{elk} punt $O$ geldt
        $\vect{AB} = \vect{OB} - \vect{OA}$.
  \item Bewijs de twee karakteriseringen van het \omterm{prop:g10:coordgeom:midpoint}{midden} $I$ van $[AB]$:
        \[
        \vect{IA} + \vect{IB} = \vec 0
        \qquad\text{en}\qquad
        \vect{OI} = \tfrac12\left(\vect{OA} + \vect{OB}\right)
        \ \text{voor elke } O .
        \]
  \item Zoek met $\vect{AB} = \vect{DC}$ het vierde hoekpunt $D$ van het
        parallellogram $ABCD$ terug, met $A(-1, 2)$, $B(3, 4)$, $C(6, 0)$ ---
        en vergelijk met het antwoord via de middentruc uit
        \cref{pb:g10:coordgeom:1}.
  \item Een vaste \omterm{def:g10:vectors:vector}{vector} $\vec u$ bij elk punt van het vlak optellen voert
        welke transformatie uit --- dezelfde die met twee spiegels opdook in
        de weekendopgave over de twee spiegels en met twee halve draaiingen
        in de weekendopgave over halve draaiingen, beide in het
        onderbouwvolume? Bereken het \omterm{def:g10:functions:function}{beeld} van $(1, 2)$ onder
        $\vec u = (3, -1)$.
\end{enumerate}

\medskip
\noindent\textbf{Deel II --- Het zwaartepunt, derde bewijs.}
Definieer het punt $G$ van een driehoek $ABC$ door de sierlijke voorwaarde
\[
\vect{GA} + \vect{GB} + \vect{GC} = \vec 0 .
\]
\begin{enumerate}[resume]
  \item Toon met de oorsprongtruc aan dat die voorwaarde $G$ ondubbelzinnig
        vastlegt: $\vect{OG} = \frac13\left(\vect{OA} + \vect{OB} +
        \vect{OC}\right)$ voor elke $O$ --- zodat de \omterm{def:g10:coordgeom:system}{coördinaten} van $G$ de
        \emph{gemiddelden} van die van de hoekpunten zijn (de weekendopgave
        over het zwaartepunt in het onderbouwvolume zag het al met getallen).
  \item Bewijs $\vect{AG} = \frac13\left(\vect{AB} + \vect{AC}\right)$, en,
        met $K$ voor het \omterm{prop:g10:coordgeom:midpoint}{midden} van $[BC]$, $\vect{AK} =
        \frac12\left(\vect{AB} + \vect{AC}\right)$. Besluit:
        $\vect{AG} = \frac23\,\vect{AK}$ --- de tweederdestelling, opnieuw
        bewezen.
  \item Waarom gaan de zwaartelijnen uit $B$ en uit $C$ door \emph{dezelfde}
        $G$, zonder ook maar één berekening extra? Vergelijk de drie
        beschikbare bewijzen van deze stelling --- het verborgen
        parallellogram uit jaar 8, de \omterm{def:g10:coordgeom:system}{coördinaten} en de \omterm{def:g10:vectors:vector}{vectoren} --- in telkens
        één zin.
  \item Numerieke controle: $A(1, 1)$, $B(5, 2)$, $C(3, 6)$. Bereken $G$,
        daarna $K$, en ga $\vect{AG} = \frac23 \vect{AK}$ na met het
        criterium voor collineariteit (\cref{thm:g10:vectors:det}).
  \item De natuurkunde leest $G$ als het evenwichtspunt van drie gelijke
        massa's in de hoekpunten. Verdubbel de massa in $A$ (massa's
        $2, 1, 1$): het evenwichtspunt wordt
        $\vect{OG'} = \frac{2\vect{OA} + \vect{OB} + \vect{OC}}{4}$. Bereken
        $G'$ voor de driehoek van vraag 9 en vergelijk zijn ligging met die
        van $G$.
\end{enumerate}

\medskip
\noindent\textbf{Deel III --- Collineariteit aan het werk.}
\begin{enumerate}[resume]
  \item Zijn $\vec u(3, -2)$ en $\vec v(-6, 4)$ \omterm{def:g10:vectors:collinear}{collineair}? En $\vec w(2, 5)$
        en $\vec z(4, 9)$? Beslis met het determinantcriterium.
  \item Liggen de punten $A(1, 2)$, $B(3, 5)$, $C(7, 11)$ op één rechte?
  \item Voorbij Varignon (\cref{exo:g10:vectors:10}): in elke vierhoek
        $ABCD$ verbinden de \emph{bimedianen} de \omterm{prop:g10:coordgeom:midpoint}{middens} van overstaande
        zijden ($[IK]$ en $[JL]$). Toon met plaatsvectoren
        ($\vect{OI} = \frac12(\vect{OA} + \vect{OB})$ enzovoort) aan dat
        beide bimedianen hetzelfde \omterm{prop:g10:coordgeom:midpoint}{midden} hebben, met plaatsvector
        $\frac14\left(\vect{OA} + \vect{OB} + \vect{OC} +
        \vect{OD}\right)$: het zwaartepunt van de vierhoek.
  \item Thales in één regel: $M$ en $N$ zijn de punten bepaald door
        $\vect{AM} = \frac13\vect{AB}$ en $\vect{AN} = \frac13\vect{AC}$.
        Bereken $\vect{MN}$ in termen van $\vect{BC}$ en besluit tot
        evenwijdigheid en verhouding.
  \item Veralgemeen vraag 14 tot een willekeurige verhouding $k$, en zeg in
        één zin hoe het vectorrekenen de middenparallelstelling en de stelling
        van Thales in één klap opslokt.
\end{enumerate}

\medskip
\noindent\textbf{Deel IV --- Krachten en snelheden.}
\begin{enumerate}[resume]
  \item Een boot wordt door een stroming van $3$ km/u pal naar het oosten
        geduwd terwijl er met $4$ km/u pal naar het noorden geroeid wordt.
        Geef de resulterende snelheidsvector, zijn grootte, en beschrijf de
        werkelijke koers. (Er duikt een zeer oud drietal op.)
  \item Op een ring werken drie krachten: $\vec u(2, 1)$, $\vec v(-3, 2)$,
        $\vec w(1, -3)$. Bereken de resultante. Wat doet de ring? En waarom
        is de \omterm{def:g10:vectors:sum}{som} van de pijlen langs elke \emph{gesloten} veelhoek altijd
        $\vec 0$ (Chasles)?
  \item Twee krachten $\vec u(4, -1)$ en $\vec v(-1, 3)$ werken op een haak.
        Welke derde kracht houdt de haak in evenwicht?
  \item Een zwemmer steekt een rivier over en zwemt recht naar de overkant
        met $2$ m/s; de stroming loopt met $1.5$ m/s. Geef de resulterende
        snelheid en haar grootte. De rivier is $50$ m breed: hoe lang duurt
        de overtocht, en hoever stroomafwaarts komt de zwemmer aan land?
  \item Slotstuk: één voorwerp, drie kostuums --- verplaatsingspijl,
        coördinatenpaar, kracht of snelheid. Noem de drie klassieke
        stellingen die deze opgave met vectoralgebra opnieuw bewees, en noem
        de ene meetkundige grootheid die de \omterm{def:g10:vectors:vector}{vectoren} van dit hoofdstuk nog
        niet kunnen meten --- het gereedschap dat dat verhelpt komt met het
        scalair product, in jaar 11.
\end{enumerate}
\end{problem}
